\section{理论基础}
\label{ch:theory}

\subsection{进化计算与LLM的结合}

LLM改变了进化计算的两个关键环节：LLM是强语义变异器（根据失败原因生成有针对性的修改），LLM可以参与评价与解释（生成评价、原则、批评）。

% ====================================================================
% 图2.1: Agent Evolution 通用进化闭环
% ====================================================================
\begin{figure}[htbp]
\centering
\begin{tikzpicture}[
    block/.style={rectangle, rounded corners=6pt, draw=#1!70!black, fill=#1!12,
        text width=3cm, minimum height=1.2cm, align=center, font=\small\bfseries},
    arr/.style={-{Stealth[length=3mm]}, thick, color=gray!70},
    label/.style={font=\scriptsize, color=gray!60, align=center},
]
    % Center loop
    \node[block=cat1] (gen) at (0,0) {生成候选\\Generate};
    \node[block=cat2] (eval) at (5,0) {评估\\Evaluate};
    \node[block=cat3] (sel) at (5,-3.5) {选择/保留\\Select \& Archive};
    \node[block=cat4] (mut) at (0,-3.5) {变异/更新\\Mutate \& Update};

    % Loop arrows
    \draw[arr] (gen.east) -- (eval.west) node[midway, above, label] {执行轨迹 $x_t$};
    \draw[arr] (eval.south) -- (sel.north) node[midway, right, label] {反馈 $y_t = V(x_t)$};
    \draw[arr] (sel.west) -- (mut.east) node[midway, below, label] {Archive $A_{t+1}$};
    \draw[arr] (mut.north) -- (gen.south) node[midway, left, label] {更新 $z_{t+1}$};

    % External inputs
    \node[block=gray, text width=2.5cm] (task) at (-4,0) {任务分布\\$P(\tau)$};
    \node[block=gray, text width=2.5cm] (safety) at (9,-3.5) {安全约束\\$C(x_t, z_t)$};

    \draw[arr, dashed, color=cat1!50] (task.east) -- (gen.west);
    \draw[arr, dashed, color=cat4!50] (safety.west) -- (sel.east);

    % Inner details
    \node[font=\scriptsize, color=cat1!50!black, text width=2.8cm, align=left] at (0,-1.5)
        {LLM Rewrite\\MCTS\\RL Update\\Code Patch};
    \node[font=\scriptsize, color=cat2!50!black, text width=2.8cm, align=left] at (5,-1.5)
        {Benchmark\\Unit Test\\Judge\\Human};

\end{tikzpicture}
\caption{Agent Evolution通用进化闭环。状态$z=(\theta, c, g, m, A)$在``生成—评估—选择—变异''循环中迭代更新，受任务分布$P(\tau)$和安全约束$C$的联合调节。}
\label{fig:evolution-loop}
\end{figure}

\subsection{自我指涉与G\"{o}del机器}

现代LLM agent与原始G\"{o}del机器的差距在于用经验评估替代形式证明。G\"{o}del Agent、DGM和SICA实现了``实用G\"{o}del化''：不要求全局最优证明，只要求每次变体在隔离评估中相对可靠地改进。

\subsection{元学习与自我改进的理论框架}

LLM agent可表示为三层策略：底层$\theta$（模型参数）、中层$\varphi$（提示词/工具/代码）、上层$m$（记忆/技能/原则）。广义策略$\pi(a|s; \theta, \varphi, m)$的自演化即更新$(\theta, \varphi, m)$的子集。

\subsection{自进化的数学形式化}

智能体系统状态$z=(\theta, c, g, m, A)$，更新过程：
\begin{equation}
z_{t+1}, A_{t+1} = S\Big(A_t \cup \{U_k(z_t, x_t, y_t)\}_{k=1}^K;\ V, C, D\Big)
\end{equation}
其中$U_k$为变异器，$S$为选择器，$V$为评估器，$C$为安全约束，$D$为多样性约束。
